In this paper, we set out to answer the question of: "how can a conceptual model of the Linux operating system be constructed that is detailed enough to represent and analyse the interactions between competing AI agents through filesystem operations and process management?" We addressed this across three main bodies of work: a survey of the Linux commandline and filesystem as the field of operations, the construction of the OSM and AESM as analytical frameworks, and the derivation of a set of theoretical strategies and limitations from those models.

The commandline and filesystem sections gave us the framework for the model. The Linux filesystem, with its directory hierarchy, inode-based metadata, VFS abstraction layer, and pseudo-filesystems such as \texttt{/proc}, provides a rich and well-documented surface through which processes observe, modify, and contest the state of a system. The commandline survey uncovered some specific operations available to an agent running locally, from process monitoring and scheduling to file manipulation and permission management, grounding the model in concrete and verifiable capabilities rather than abstract assumptions.

We then formalised this baseline in the OSM as a memory-state machine, representing filesystem state as a grid of labelled block categories and process activity as a set of transitions between those states. The AESM then introduced two competing agents into this framework, from which we derived the criteria for being alive or dead in a computational sense: an agent is alive if it is executing unaltered code, retains access to the filesystem, and has not had its main processes terminated. From these criteria, and from the properties of the underlying system, a set of practical strategies emerged: watchdog processes for process resilience, distributed backups for file integrity, dummy files and processes for misdirection, and map poisoning as an offensive counterpart to state map monitoring. The arms race dynamic that arises from these strategies, and the resource-efficiency hypothesis it implies, are what we consider the primary theoretical contribution of this paper.

The discussion surfaced a set of honest limitations. The model is a static snapshot rather than a dynamic process, treats agent state as binary rather than graduated, omits RAM as a strategic medium, leaves privilege level unspecified, and assumes a symmetry between the two agents that is unlikely to reflect real deployments. None of these limitations undermine what we found, but they do bound how confidently we can transfer those conclusions to a live system. They also define the most productive directions for future work: extending the model with time series, intermediate states, and memory representation; specifying the privilege context; and eventually validating the derived strategies in a simulation. Considerations that also underline the importance of an actual virtual machine simulation as one of if not the most important next steps for this research.

What emerges most clearly from all of this is that the Linux operating system, taken as a whole, turns out to be a surprisingly well-suited arena for AI interaction. The filesystem provides a persistent, structured record of activity that any agent can read and exploit; the process table provides a live map of active computation that can be monitored, disrupted, or obscured; and the kernel's own accounting mechanisms, inode timestamps, scheduling statistics, and the \texttt{/proc} interface, give both agents access to the same ground truth about system state. This means that neither agent has a structural information advantage, so the outcome ultimately comes down to the quality of each agent's strategy rather than its access to privileged data.

This has a broader implication for the field of autonomous cybersecurity. As AI capabilities continue to develop along the trajectory documented in the introduction, the scenarios described here will not remain theoretical for long. Designing defensive systems that are robust against an adversary capable of the strategies we have derived here requires exactly the kind of detailed, operation-level model that this work has begun to build. The OSM and AESM, even in their current simplified form, give us a concrete starting point for that design work, and the limitations we identified in the discussion give us a clear roadmap for the iterations that should follow.